\chapter{Dry Eyes}
\index{Dry Eyes}

\section*{Definition and Overview}
Dry eye disease is a common condition in which the eyes do not produce enough tears, or the tears evaporate too quickly, leaving the surface of the eye inadequately lubricated. The tear film has three layers: an oily outer layer that prevents evaporation, a watery middle layer that provides moisture, and a mucin layer that helps tears spread evenly. When any of these layers is deficient, the eye becomes irritated, gritty, and uncomfortable. Dry eye is usually a chronic condition, but it is very treatable, and most people can achieve substantial relief with simple measures and appropriate treatment.

\section*{Epidemiology}
Dry eye is one of the most common reasons for visiting an eye care professional, affecting a large proportion of adults, with prevalence increasing with age. It is more common in women, particularly after menopause, and is associated with ageing, screen use, dry environments, and certain medications. The condition is often underestimated because symptoms range from mild irritation to significant pain and visual disturbance, and many people self-treat for years without seeking professional care.

\section*{Aetiology and Pathophysiology}
Dry eye occurs when there is a mismatch between tear production and tear loss. It is divided into two main types, which often coexist.
\begin{itemize}
    \item \textbf{Aqueous-deficient dry eye:} the lacrimal glands do not produce enough watery tears, as occurs in Sjögren's syndrome, after ageing, or with some medications
    \item \textbf{Evaporative dry eye:} the oily layer produced by the meibomian glands is deficient, so tears evaporate too quickly; this is the most common type and is linked to blepharitis and meibomian gland dysfunction
    \item \textbf{Ageing:} tear production naturally declines with advancing years
    \item \textbf{Hormonal changes:} oestrogen and androgen levels influence tear glands, explaining the higher rates in women after menopause
    \item \textbf{Environmental factors:} air conditioning, wind, smoke, dry climates, and prolonged screen use increase evaporation
    \item \textbf{Medications:} antihistamines, decongestants, antidepressants, and some blood pressure medicines reduce tear production
    \item \textbf{Medical conditions:} rheumatoid arthritis, Sjögren's syndrome, thyroid disease, and diabetes are associated with dry eye
    \item \textbf{Contact lens wear:} lenses absorb tears and irritate the ocular surface
\end{itemize}

\section*{Clinical Features}
Symptoms may be constant or fluctuate with the environment and time of day.
\begin{itemize}
    \item Grittiness, burning, or stinging in the eyes
    \item A feeling of something in the eye, as if there is sand or grit
    \item Redness, itching, and irritation
    \item Blurred vision that improves with blinking
    \item Eye fatigue, especially after reading or screen work
    \item Paradoxical watery eyes: reflex tearing in response to dryness
    \item Discomfort when wearing contact lenses
    \item Sensitivity to light
\end{itemize}

\section*{Differential Diagnosis}
Several other conditions cause similar symptoms and may coexist with dry eye.
\begin{itemize}
    \item Allergic conjunctivitis, which causes itching, redness, and watering
    \item Infective conjunctivitis, with discharge and crusting
    \item Blepharitis, inflammation of the eyelid margins
    \item Contact lens-related problems, including corneal irritation and infection
    \item Corneal abrasions or recurrent corneal erosion
    \item Refractive error, causing eye strain and fatigue rather than dryness
\end{itemize}

\section*{When to Seek Medical Care}
See a GP, optometrist, or ophthalmologist if dry eye symptoms are persistent, worsening, or interfering with daily life, or if self-care measures are not helping. A professional can identify the underlying cause and recommend appropriate treatment.

\textbf{Contact your local emergency services for:}
\begin{itemize}
    \item Sudden severe eye pain
    \item Sudden loss of vision or major change in vision
    \item A red eye with discharge, which may indicate a serious infection
    \item An injury or chemical in the eye
\end{itemize}

\section*{Diagnosis and Investigations}
Diagnosis of dry eye is usually made by an eye care professional using a combination of symptom assessment and examination.
\begin{itemize}
    \item \textbf{History:} symptoms, medication use, screen time, contact lens wear, and any underlying conditions
    \item \textbf{Slit-lamp examination:} magnified inspection of the eye surface, tear film, and eyelid margins
    \item \textbf{Tear film tests:} measurements of tear production, such as the Schirmer test, and of tear break-up time to assess evaporation
    \item \textbf{Staining:} dyes that reveal damage to the surface of the cornea and conjunctiva
    \item \textbf{Meibomian gland assessment:} evaluation of the oil-producing glands for dysfunction
    \item \textbf{Blood tests:} where an underlying condition such as Sjögren's syndrome or thyroid disease is suspected
\end{itemize}

\section*{Management}
Treatment is stepped and tailored to the severity and type of dry eye.
\begin{itemize}
    \item \textbf{Artificial tears:} lubricating eye drops used regularly, or as needed; preservative-free drops are preferable for frequent use
    \item \textbf{Lubricating gels and ointments:} thicker preparations used at night for prolonged coverage
    \item \textbf{Eyelid hygiene:} warm compresses and gentle lid massage improve meibomian gland function in evaporative dry eye
    \item \textbf{Environmental measures:} humidifiers, taking screen breaks, and avoiding direct air flow from fans and air conditioners
    \item \textbf{Prescription eye drops:} anti-inflammatory drops, including cyclosporine and lifitegrast, reduce ocular surface inflammation
    \item \textbf{Punctal plugs:} tiny plugs placed in the tear drainage ducts to keep tears on the eye longer
    \item \textbf{Managing contributing factors:} reviewing medications, optimising contact lens wear, and treating blepharitis or underlying conditions
    \item \textbf{Advanced treatments:} intense pulsed light therapy and other in-clinic procedures for meibomian gland disease
\end{itemize}

\section*{Complications}
\begin{itemize}
    \item Corneal damage, including superficial punctate keratopathy and, in severe cases, corneal ulceration or scarring
    \item Increased risk of eye infections because the protective tear film is compromised
    \item Reduced quality of life, with difficulty reading, driving, and working
    \item Impaired vision from chronic surface damage
    \item Psychological effects of chronic discomfort
\end{itemize}

\section*{Prognosis and Outcomes}
Dry eye is a chronic condition for most people, but with appropriate treatment the vast majority achieve good control of their symptoms. The condition rarely threatens sight in properly managed patients, although severe cases require ongoing specialist care. Treatment is usually long-term, and symptoms may flare with changes in environment, season, medication, or health. Early diagnosis and consistent self-care offer the best chance of comfortable eyes.

\section*{Prevention and Self-Care}
\begin{itemize}
    \item Follow the 20-20-20 rule: every 20 minutes of screen time, look at something 20 feet away for 20 seconds
    \item Blink regularly and fully, especially during screen work
    \item Use a humidifier in dry or air-conditioned rooms
    \item Protect eyes from wind, sun, and smoke with wrap-around sunglasses
    \item Stay well hydrated and consider omega-3 supplements, which may help some people
    \item Perform regular warm-compress eyelid hygiene
    \item Have regular eye examinations, particularly over the age of 40
\end{itemize}

\section*{Sources and Further Reading}
The following reputable sources provide further information for healthcare students and patients seeking reliable, current advice about dry eyes.
\begin{itemize}
    \item NHS. \textit{Dry eye syndrome: causes, symptoms, and treatment.} https://www.nhs.uk/conditions/dry-eyes/
    \item Mayo Clinic. \textit{Dry eyes: symptoms, causes, and treatment.} https://www.mayoclinic.org/diseases-conditions/dry-eyes/
    \item MedlinePlus. \textit{Dry eye syndrome.} https://medlineplus.gov/dryeyesyndrome.html
    \item American Academy of Ophthalmology. \textit{Dry eye: patient education.} https://www.aao.org/eye-health/diseases/dry-eye
    \item MSD Manual. \textit{Dry eye: professional overview.} https://www.msdmanuals.com/
    \item Tear Film \& Ocular Surface Society. \textit{International dry eye workshop information.} https://www.tearfilm.org/
\end{itemize}
